\subsection{Beliefs and trading}\label{app:beliefs}

\Urlmuskip=0mu plus 2mu
\setlength{\emergencystretch}{1.5em}

\paragraph{Beliefs, prices and trading.}\label{par:a2-scope}
The belief lemma holds for every order profile, in or out of equilibrium; the price and trading lemmas hold
for every tuple that meets their stated conditions.

\paragraph{Conventions for Appendices A.2 to A.9.}
The conditions (E1) to (E6) are those of Definition~\ref{def:equilibrium}.
\begin{itemize}[leftmargin=2em]
\item A \emph{profile} is a pair \((\sigma_H,\sigma_L)\) of order distributions, pure or mixed, with wrong-signed
orders allowed. A \emph{candidate} is a tuple \(\mathcal E\) that satisfies (E1) to (E4) and (E6) for some
profile. An equilibrium is a candidate that also satisfies (E5).
\item A deviation by the investor holds the candidate's price function \(P\) and preparation rule \(\zeta\) fixed
and changes only the order that reaches them; a comparison across incumbents solves the equilibrium in each
economy to compare outcomes across economies.
\item A uniqueness proof derives its bounds for every candidate. It does not use the proposed tuple.
\item A constructed tuple satisfies each equilibrium condition by direct substitution, and its proof cites no lemma whose premises include the condition at issue.
\item Write \(w_j=t_j-t_0\) for the gain in expected proceeds that entry brings in class \(j\), and
\(w(\mu)=(1-\mu)w_L+\mu w_H=t_L-t_0+\DT\mu\) for its average at the belief \(\mu\). Proposition~\ref{prop:returns}
and Section~\ref{sec:profits} give \(w_H\ge w_L\ge0\).
\end{itemize}

% STRUCTURE: the symbols delta_q and b_0 and the odds form eq:odds <- Section 3.2
\paragraph{Orders and posteriors.}
Write \(\delta_q\) for the point mass at the order \(q\), and \(b_0=(1-\lambda)(1-2\beta)\) for the probability
of a liquidity abstention. Under any order distributions \(\sigma_H,\sigma_L\), Bayes' rule in odds form reads
\begin{equation}\label{eq:odds}
\frac{\mu(x)}{1-\mu(x)}=\frac{\pi}{1-\pi}\cdot\frac{f_H(x)}{f_L(x)}.
\end{equation}

% STRUCTURE: Lemma 1 <- Section 3.2
\begin{lemma}[Order probabilities and beliefs]\label{lem:beliefs}
Fix any order distributions \(\sigma_H\) and \(\sigma_L\).
\begin{enumerate}[label=(\alph*),leftmargin=2em]
\item \emph{Full support.} \(f_j(x)\ge(1-\lambda)\eta(x)>0\) for every order \(x\) and class \(j\).
\item \emph{Nonzero orders.} For \(x\in\{-1,1\}\), \(b\le f_j(x)\le b+\lambda\). Hence
\(m(\pi)\le\mu(x)\le M(\pi)\).
\item \emph{The zero order.} \(b_0\le f_j(0)\le b_0+\lambda\). If \(\beta\le1/3\), then \(b_0\ge b\), hence
\(m(\pi)\le\mu(0)\le M(\pi)\).
\item \emph{Full orders.} If \(\sigma_H=\delta_1\) and \(\sigma_L=\delta_{-1}\), the posteriors after the orders
\((-1,0,1)\) are \(\bigl(m(\pi),\pi,M(\pi)\bigr)\).
\end{enumerate}
\end{lemma}

\begin{proof}[Proof of Lemma~\ref{lem:beliefs}]
Fix any profile \((\sigma_H,\sigma_L)\). By \eqref{eq:odds}, the posterior \(\mu(x)\) increases strictly in the
likelihood ratio \(f_H(x)/f_L(x)\). The odds of \(M(\pi)\) are \(\pi(1+u)/(1-\pi)\), and the odds of \(m(\pi)\) are
\(\pi/[(1-\pi)(1+u)]\).

\emph{Part (a).} The liquidity trader fills the slot with probability \(1-\lambda>0\). It places the order \(x\)
with probability \(\eta(x)\), and \(\eta(x)>0\) for every \(x\) because \(0<\beta<1/2\). The informed term
\(\lambda\sigma_j(x)\) is nonnegative. Hence \(f_j(x)\ge(1-\lambda)\eta(x)>0\) for every order and class. Every
order has positive probability in both classes, so (E1) defines \(\mu(x)\) at every order, and no off-path belief
enters any condition. Also \(0<\mu(x)<1\), since \(0<\pi<1\).

\emph{Part (b).} For \(x=\pm1\), \((1-\lambda)\eta(x)=b\) and \(0\le\lambda\sigma_j(x)\le\lambda\). Hence
\(b\le f_j(x)\le b+\lambda\), and
\[
\frac1{1+u}=\frac{b}{b+\lambda}\le\frac{f_H(x)}{f_L(x)}\le\frac{b+\lambda}{b}=1+u .
\]
By \eqref{eq:odds}, \(m(\pi)\le\mu(x)\le M(\pi)\). This step uses only \(0<\beta<1/2\).

\emph{Part (c).} For \(x=0\), \((1-\lambda)\eta(0)=b_0\), so \(b_0\le f_j(0)\le b_0+\lambda\), and the likelihood
ratio lies in \([(1+\lambda/b_0)^{-1},\,1+\lambda/b_0]\). The inequality \(b_0\ge b\) is \(1-2\beta\ge\beta\), that
is \(\beta\le1/3\). Then \(\lambda/b_0\le\lambda/b=u\), and the argument of part (b) gives
\(m(\pi)\le\mu(0)\le M(\pi)\). If \(\beta>1/3\), then \(b_0<b\). The profile \(\sigma_H=\delta_0\) with
\(\sigma_L(0)=0\) gives \(f_H(0)=b_0+\lambda\) and \(f_L(0)=b_0\). The ratio is \(1+\lambda/b_0>1+u\), so
\(\mu(0)>M(\pi)\), and the bound fails.

\emph{Part (d).} The profile \(\sigma_H=\delta_1\), \(\sigma_L=\delta_{-1}\) gives \(f_H=(b,\,b_0,\,b+\lambda)\) and
\(f_L=(b+\lambda,\,b_0,\,b)\) on the orders \((-1,0,1)\). The likelihood ratios are \((1+u)^{-1}\), \(1\) and
\(1+u\), so \((\mu(-1),\mu(0),\mu(1))=(m(\pi),\pi,M(\pi))\).

\emph{Beliefs at a price.} Fix a price \(\bar P\) in the image of \(P\), and let \(X(\bar P)=\{x:P(x)=\bar P\}\).
By (E2), the challenger's belief \(\bar\mu(\bar P)\) has the odds
\[
\frac{\bar\mu(\bar P)}{1-\bar\mu(\bar P)}=\frac{\pi}{1-\pi}\cdot
\frac{\sum_{x\in X(\bar P)}f_H(x)}{\sum_{x\in X(\bar P)}f_L(x)} .
\]
A ratio of two sums of positive terms lies between the smallest and the largest ratio of its terms. Hence
\(\bar\mu(\bar P)\) lies between the smallest and the largest \(\mu(x)\) over \(X(\bar P)\). Every bound of parts (b)
and (c) therefore holds for \(\bar\mu(\bar P)\).

Only part (c) uses \(\beta\le1/3\). Parts (a), (b) and (d) hold for every \(\beta<1/2\).
\end{proof}

% STRUCTURE: Corollary A.1 <- Section 3.2
\begin{corollary}[High-cost preparation under every profile]\label{cor:exclusion}
Let \(0<\beta\le1/3\) and \(g_H\ge g_L\). If \(B(M(\pi))<c_H\), then no profile that satisfies (E1)--(E3), in or out of equilibrium,
has high-cost preparation after any order.
\end{corollary}

\begin{proof}[Proof of Corollary~\ref{cor:exclusion}]
Let \(\beta\le1/3\), \(g_H\ge g_L\) and \(B(M(\pi))<c_H\). Fix any profile and any tuple that satisfies (E1) to
(E3). By Lemma~\ref{lem:beliefs}(b)--(c), \(\mu(x)\le M(\pi)\) at every order, and by the step on beliefs at a price,
\(\bar\mu(\bar P)\le M(\pi)\) at every price. Since \(g_H\ge g_L\), \(B(\mu)=g_L+\mu(g_H-g_L)\) is nondecreasing in
\(\mu\). Hence \(B(\bar\mu(\bar P))\le B(M(\pi))<c_H\) at every price, and (E3) gives \(\zeta(\bar P)=0\) with no
tie. The conclusion holds under both challenger rules. The bound at the zero order is the only use of
\(\beta\le1/3\) in this proof.
\end{proof}

% STRUCTURE: Lemma 2 and its arbitrary-preparation clause <- Section 3.2
\begin{lemma}[Entry and prices]\label{lem:prices}
Let \(\DT>0\) and \(g_H>g_L\). Fix a tuple that satisfies (E1)--(E4) and (E6).
\begin{enumerate}[label=(\alph*),leftmargin=2em]
\item \emph{Floor.} \(e(x)\ge\rho\) for every order \(x\).
\item \emph{Price.} \(P(x)=t_0+e(x)\bigl[t_L-t_0+\DT\,\mu(x)\bigr]\) for every order \(x\).
\item \emph{Sufficiency.} If \(P(x)=P(x')\), then \(\mu(x)=\mu(x')\). The challenger's belief at the price \(P(x)\)
is \(\mu(x)\).
\item \emph{Entry.} \(e(x)=\rho\) if \(\mu(x)<\tau\), and \(e(x)=1\) if \(\mu(x)\ge\tau\).
\item \emph{Monotonicity.} If \(\mu(x)>\mu(x')\), then \(e(x)\ge e(x')\) and
\(P(x)-P(x')\ge\rho\DT\,[\mu(x)-\mu(x')]>0\).
\end{enumerate}
\end{lemma}
When an indifferent high-cost challenger may prepare with any probability, part~(d) at \(\mu(x)=\tau\) allows any \(e(x)\in[\rho,1]\); the other parts hold under both
challenger rules.

\begin{proof}[Proof of Lemma~\ref{lem:prices}]
Fix a candidate, and let \(\DT>0\) and \(g_H>g_L\). The proof has a preliminary step and five parts.

\emph{Step 0: the terminal payoff given the class and the order.} The challenger's choice depends only on the
price \(P(x)\) and its cost \(C\). The cost is independent of the class, the values, the incumbent's value and the
slot draws. Given the class \(j\) and the order \(x\), the challenger therefore prepares with probability
\(e(x)=\rho+(1-\rho)\zeta(P(x))\), and this event is independent of \((\theta,R)\). The order depends on the class
and the slot draws only, so it carries no information about \((\theta,R)\) given \(j\). Bids are truthful by (E6).
With entry, expected proceeds in class \(j\) are \(t_j\); without entry, they are \(t_0\). Hence
\begin{equation}\label{eq:a2-terminal}
\E[V_T\mid j,X=x]=t_0+e(x)\,w_j .
\end{equation}

\emph{Part (a): the floor.} The zero-cost challenger prepares after every price (Definition~\ref{def:equilibrium}),
and it has probability \(\rho\). Hence \(e(x)\ge\rho\) at every order.

\emph{Part (b): the price.} Average \eqref{eq:a2-terminal} over the class with the weights \(\mu(x)\) and
\(1-\mu(x)\) of (E1). Condition (E4) gives \(P(x)=t_0+e(x)\,w(\mu(x))\), which is \eqref{eq:price}. This step uses
only (E1), (E4), the truthful bids (E6) of Step 0, and the independence of the cost.

\emph{Part (c): sufficiency.} Let \(P(x)=P(x')\). Entry depends on the order only through the price, so
\(e(x)=e(x')=:\bar e\), and \(\bar e\ge\rho>0\) by part (a). Part (b) gives \(\bar e\,\DT\,[\mu(x)-\mu(x')]=0\), so
\(\mu(x)=\mu(x')\). By the step on beliefs at a price in the proof of Lemma~\ref{lem:beliefs}, \(\bar\mu(P(x))\)
lies between the smallest and the largest posterior of the orders with the price \(P(x)\). These posteriors are
equal, so \(\bar\mu(P(x))=\mu(x)\). The price therefore reveals the posterior.

\emph{Part (d): entry.} Since \(g_H>g_L\), \(B(\mu)>c_H\) exactly when \(\mu>\tau\). By part (c), the
challenger's gross profit at \(P(x)\) is \(B(\mu(x))\). Condition (E3) gives \(e(x)=\rho\) if \(\mu(x)<\tau\) and
\(e(x)=1\) if \(\mu(x)>\tau\). At \(\mu(x)=\tau\), \(e(x)=1\) when an indifferent high-cost challenger prepares, and \(e(x)\) can take any value in
\([\rho,1]\) when an indifferent high-cost challenger may prepare with any probability.

\emph{Part (e): monotonicity.} Let \(\mu(x)>\mu(x')\). If \(\mu(x')<\tau\), then \(e(x')=\rho\le e(x)\). If
\(\mu(x')\ge\tau\), then \(\mu(x)>\tau\) and \(e(x)=1\ge e(x')\). Part (b) gives
\begin{equation}\label{eq:a2-increment}
P(x)-P(x')=e(x)\,\DT\,[\mu(x)-\mu(x')]+[e(x)-e(x')]\,w(\mu(x')) .
\end{equation}
The second term is nonnegative, because \(w(\mu)\ge w_L\ge0\). The first term is at least
\(\rho\DT[\mu(x)-\mu(x')]>0\). Prices therefore rise strictly with the posterior.
\end{proof}

\paragraph{Constructed tuples.}\label{par:constructed-tuples}
Let \(\DT>0\). Take a profile, its Bayes posteriors \(\mu(x)\), and entry \(e(x)\in[\rho,1]\) such that
\(e(x)\ge e(x')\) whenever \(\mu(x)>\mu(x')\). Set \(P(x)=t_0+e(x)[w_L+\DT\mu(x)]\), where \(w_L=t_L-t_0\ge0\).
Then prices rise strictly with the posterior, and
\[
P(x)=P(x')\quad\Longleftrightarrow\quad
\mu(x)=\mu(x')\ \text{and}\ e(x)=e(x').
\]
The preparation rule \(\zeta(P(x))=[e(x)-\rho]/(1-\rho)\) is a function of the price, and the belief at \(P(x)\)
is \(\mu(x)\). If also \(g_H>g_L\), threshold-consistent entry is a function of the posterior when an indifferent
high-cost challenger prepares. When an indifferent high-cost challenger may prepare with any probability,
orders at \(\tau\) with different entry have different prices, each revealing \(\tau\).

\begin{proof}
Let \(\mu(x)>\mu(x')\). Then \(e(x)\ge e(x')\), and \eqref{eq:a2-increment} gives
\(P(x)-P(x')\ge e(x)\DT[\mu(x)-\mu(x')]\ge\rho\DT[\mu(x)-\mu(x')]>0\), because \(w(\mu)\ge w_L\ge0\). Let
\(\mu(x)=\mu(x')\). Then \(P(x)-P(x')=[e(x)-e(x')]\,w(\mu(x))\), and \(w(\mu(x))=w_L+\DT\mu(x)>0\), because
\(\mu(x)>0\) by Lemma~\ref{lem:beliefs}(a). Equal prices therefore force equal posteriors and equal entry, and the
price formula gives the converse. So \(\zeta\) is a function of the price, and \(e(x)\in[\rho,1]\) puts its values in
\([0,1]\). At a price \(\bar P\), all orders with the price \(\bar P\) share one posterior, so by \eqref{eq:odds}
they share one likelihood ratio \(f_H(x)/f_L(x)\). The ratio of sums in the step on beliefs at a price, in the
proof of Lemma~\ref{lem:beliefs}, equals that likelihood ratio, and (E2) gives the common posterior. Under the
first challenger rule, threshold-consistent entry is \(\rho\) below \(\tau\) and one at and above \(\tau\), a function
of the posterior. Under the second,
the identity at equal posteriors gives distinct prices to distinct entries at \(\tau\), and every order with one of
these prices has the posterior \(\tau\).
\end{proof}

% STRUCTURE: Lemma 3 <- Section 3.3
\begin{lemma}[Informed trading]\label{lem:trading}
Let \(\DT\ge0\). Fix a tuple that satisfies (E1), (E4) and (E6), with entry \(e(x)\in[\rho,1]\) after every order.
Conditional on its arrival, the investor in each class earns
\[
\begin{aligned}
U_H(1)&=e(1)\DT[1-\mu(1)]-k,&\quad U_H(0)&=0,&\quad U_H(-1)&=-e(-1)\DT[1-\mu(-1)]-k,\\
U_L(-1)&=e(-1)\DT\,\mu(-1)-k,&\quad U_L(0)&=0,&\quad U_L(1)&=-e(1)\DT\,\mu(1)-k.
\end{aligned}
\]
Write \(K_H=e(1)\DT[1-\mu(1)]\) and \(K_L=e(-1)\DT\,\mu(-1)\) for the gross returns of the two correct orders.
\begin{enumerate}[label=(\alph*),leftmargin=2em]
\item A wrong-signed order earns at most \(-k\), so abstention strictly beats it.
\item Class \(H\) buys if \(K_H>k\) and abstains if \(K_H<k\). At \(K_H=k\), the purchase and abstention are its
two best orders. Class \(L\) sells if \(K_L>k\) and abstains if \(K_L<k\). At \(K_L=k\), the sale and abstention
are its two best orders.
\item Under every candidate profile, \(\rho\DT[1-M(\pi)]\le K_H\le\DT\) and \(\rho\DT\,m(\pi)\le K_L\le\DT\).
Hence \(\rho\,\kappa(\pi)\DT\le K_j\le\DT\), where \(\kappa(\pi)=\min\{m(\pi),1-M(\pi)\}\).
\item \(K_H\) depends on the candidate only through \(\mu(1)\) and \(e(1)\), and \(K_L\) only through
\(\mu(-1)\) and \(e(-1)\).
\end{enumerate}
\end{lemma}

\begin{proof}[Proof of Lemma~\ref{lem:trading}]
Fix a tuple that satisfies (E1), (E4) and (E6), with \(e(x)\in[\rho,1]\) after every order, and let \(\DT\ge0\). A
deviation to the order \(q\) reaches the tuple's price \(P(q)\) and entry \(e(q)\). By \eqref{eq:a2-terminal} and
part (b) of the proof of Lemma~\ref{lem:prices},
\[
\begin{aligned}
\E[V_T\mid H,X=q]-P(q)&=e(q)\,[w_H-w(\mu(q))]=e(q)\,\DT\,[1-\mu(q)],\\
P(q)-\E[V_T\mid L,X=q]&=e(q)\,[w(\mu(q))-w_L]=e(q)\,\DT\,\mu(q).
\end{aligned}
\]
The proceeds \(t_0\) and the term \(e(q)w_L\) cancel. Multiply by \(q\) and subtract \(k|q|\). This gives the six
payoffs of the lemma, with \(K_H=e(1)\DT[1-\mu(1)]\) and \(K_L=e(-1)\DT\mu(-1)\).

\emph{Part (a).} Since \(e\ge0\), \(\DT\ge0\) and \(0\le\mu\le1\), the payoffs \(U_H(-1)=-e(-1)\DT[1-\mu(-1)]-k\)
and \(U_L(1)=-e(1)\DT\mu(1)-k\) are at most \(-k<0=U_j(0)\). No class puts weight on a wrong-signed order under
(E5).

\emph{Part (b).} By part (a), class \(H\) compares the purchase with abstention, and
\(U_H(1)-U_H(0)=K_H-k\). Class \(H\) buys with probability one if \(K_H>k\) and abstains if \(K_H<k\). If \(K_H=k\),
the purchase and abstention are its two maximizers: an indifferent investor may mix over them or abstains. Class \(L\) compares the sale with abstention, and \(U_L(-1)-U_L(0)=K_L-k\). Class \(L\) sells with
probability one if \(K_L>k\) and abstains if \(K_L<k\). If \(K_L=k\), the sale and abstention are its two maximizers:
an indifferent investor may mix over them or abstains.

\emph{Part (c).} By Lemma~\ref{lem:beliefs}(b), \(\mu(1)\le M(\pi)\) and \(\mu(-1)\ge m(\pi)\) under every profile,
for every \(\beta<1/2\). Also \(\rho\le e\le1\) and \(0\le\mu\le1\). Hence
\begin{equation}\label{eq:a2-bounds}
\rho\,\kappa(\pi)\,\DT\le\rho\,[1-M(\pi)]\,\DT\le K_H\le\DT,\qquad
\rho\,\kappa(\pi)\,\DT\le\rho\,m(\pi)\,\DT\le K_L\le\DT .
\end{equation}
The upper bounds use only \(e\le1\) and \(0\le\mu\le1\).

\emph{Part (d).} The return \(K_H\) depends on the tuple only through \((\mu(1),e(1))\), and \(K_L\) only through
\((\mu(-1),e(-1))\). Entry after the zero order enters no payoff. The ex ante payoff of the investor is
\(\lambda\,U_j(q)\), and the factor \(\lambda>0\) changes no comparison.
\end{proof}

\begin{lemma}[The auction through the spread and the threshold]\label{lem:reduction}
Fix the model's independence and information structure, a prior \(\pi\in(0,1)\), the two-point preparation cost
with \(\rho\in(0,1)\) and \(c_H>0\), a trading cost \(k>0\), and one pair of tie rules: an indifferent investor
may mix over its best orders or abstains, and an indifferent high-cost challenger prepares or may prepare with
any probability. At fixed presence, fix \(0<\lambda<1\); with research, presence is \(\hat\lambda\).
Fix \(0<\beta<1/2\). Consider two cash second-price auctions with finite-mean classes satisfying the class
ranking, continuous incumbent laws, \(\DT>0\) and \(\gamma=g_H-g_L>0\), and the same spread and preparation threshold
\(\tau=(c_H-g_L)/\gamma\). Their equilibrium sets projected onto
\((\sigma_H,\sigma_L,\mu,e)\) coincide.

The common conditions are Bayes' rule, the threshold entry rule under the chosen tie rule, and optimal orders at
\[
\begin{aligned}
U_H(1)&=e(1)\DT[1-\mu(1)]-k,& U_H(-1)&=-e(-1)\DT[1-\mu(-1)]-k,\\
U_L(-1)&=e(-1)\DT\mu(-1)-k,& U_L(1)&=-e(1)\DT\mu(1)-k,
\end{aligned}
\qquad U_j(0)=0.
\]
Each auction gives prices and price beliefs by
\[
P(x)=t_0+e(x)[w_L+\DT\mu(x)],\qquad
\zeta(P(x))=\frac{e(x)-\rho}{1-\rho},\qquad \bar\mu(P(x))=\mu(x).
\]
With research, fix also the quadratic cost coefficient \(a>0\) and cap \(0<\lb<1\), and require optimal orders
after a hypothetical success even at zero research. The equilibrium sets projected onto
\((\hat\lambda,\sigma_H,\sigma_L,\mu,e)\) coincide. The common conditions also include
\[
\hat\lambda=\min\{\lb,\Psi/a\},\qquad
\Psi=\pi\max_q U_H(q)+(1-\pi)\max_q U_L(q),
\]
with Bayes' rule evaluated at \(\hat\lambda\). These conclusions hold under all four tie rules.
\end{lemma}

\begin{proof}
For either auction, Lemma~\ref{lem:prices} makes the price posterior equal to the order posterior.
Since \(B(\mu)-c_H=\gamma(\mu-\tau)\), preparation optimality is the same threshold rule in both auctions,
including the selected action at equality. Subtracting the competitive price from the class payoff gives the
four nonzero-order payoffs in the statement by Lemma~\ref{lem:trading}. The price levels \(t_0\) and \(w_L\)
cancel. Thus every equilibrium gives a solution of the same reduced conditions.

Conversely, take a solution of those conditions. Threshold-consistent entry is nondecreasing across distinct
posteriors. The constructed-tuples paragraph supplies the displayed prices, a preparation rule that is a function
of price, and the price beliefs. In particular, orders with a common price have a common posterior and common
entry; orders at the preparation threshold with different entry have different prices. Bayes' rule at a pooled
price follows by summing order probabilities with the same likelihood ratio. Preparation is optimal because
\(\gamma>0\), competitive pricing follows from the conditional terminal payoff, and optimal orders are exactly
the reduced payoff comparisons. Truthful bidding completes the tuple. This construction works in either auction.

With research, at the candidate prices and preparation rule a deviation to intensity \(\tilde\lambda\) and optimal
success-contingent orders earns
\[
\tilde\lambda\bigl[\pi\max_qU_H(q)+(1-\pi)\max_qU_L(q)\bigr]-\frac a2\tilde\lambda^2.
\]
Its strict concavity on \([0,\lb]\) gives \(\hat\lambda=\min\{\lb,\Psi/a\}\).
At zero intensity the stipulated hypothetical-success optimality gives the same maximizers. Since the payoffs,
coefficient and cap coincide, the research condition and the price construction coincide as well.
\end{proof}

\begin{lemma}[The two numbers along a path]\label{lem:twonumbers}
Fix ranked finite-mean classes, the reserve and the two-point preparation cost. A path is a family \(F_r\) of
continuous incumbent laws indexed by strength on an interval, with \(F_{r'}\le F_r\) when \(r'>r\). Write
\(\DT(r)=\DT(F_r)\), \(\gamma(r)=g_H(F_r)-g_L(F_r)\), and, where \(\gamma(r)>0\),
\(\tau(r)=[c_H-g_L(F_r)]/\gamma(r)\). The class mean gap is fixed:
\[
\bar\Gamma=\E_H\theta-\E_L\theta=\DT(r)+\gamma(r).
\]
Each \(g_j\) and \(\gamma\) is nonincreasing, and \(\DT\) is nondecreasing. For \(r'>r\),
\[
\begin{aligned}
\DT(r')-\DT(r)&=\int_p^\infty[F_r(y)-F_{r'}(y)]D(y)\,dy,\\
g_j(F_r)-g_j(F_{r'})&=\int_p^\infty[F_r(y)-F_{r'}(y)][1-G_j(y)]\,dy.
\end{aligned}
\]
Each change is strict exactly when its integral is positive. Where \(\gamma>0\), \(\DT>0\) and \(\tau\ge0\), both
\(\tau\) and \(\DT\tau\) are nondecreasing. Their exact increments are
\[
\tau(r')-\tau(r)=
\frac{g_L(F_r)-g_L(F_{r'})+\tau(r)[\DT(r')-\DT(r)]}{\gamma(r')},
\]
\[
\DT(r')\tau(r')-\DT(r)\tau(r)
=\DT(r)[\tau(r')-\tau(r)]+\tau(r')[\DT(r')-\DT(r)].
\]
Thus the threshold rises strictly exactly when the numerator in its increment is positive; the product rises
strictly exactly when the right side of its increment is positive.

On a differentiable path the return \(J=\DT(1-\tau)\) obeys
\[
J'=\frac{\bar\Gamma(g_H-c_H)\DT'+\DT\gamma g_H'}{\gamma^2}.
\]
Its direction depends on the balance between spread growth and the decline in high-class profit. For point
classes \(h=\reg{twonum_h}\), \(\ell=\reg{twonum_ell}\), reserve \(p=\reg{twonum_p}\), high cost
\(c_H=\reg{twonum_cH}\), and \(F_r=U[r-\reg{twonum_halfwidth},r+\reg{twonum_halfwidth}]\) on
\([\reg{twonum_rlo},\reg{twonum_rhi}]\),
\[
\DT=r-\ell,\qquad g_L=0,\qquad g_H=h-r,\qquad \tau=\frac{c_H}{h-r}.
\]
Its return falls between the endpoints:
\[
\begin{aligned}
J(\reg{twonum_rlo})&=\reg{twonum_Jlo},\\
J(\reg{twonum_rhi})&=\reg{twonum_Jhi}<J(\reg{twonum_rlo}).
\end{aligned}
\]
In the uniform-incumbent example below, \(J\) rises on \([r_0,r_2]\) by the uniform case of
Proposition~\ref{prop:coexistence}(c).
\end{lemma}

\begin{proof}
The split identity \(t_j(F)+g_j(F)=\E_j\theta\) gives
\(\DT+\gamma=\bar\Gamma\). The auction integrals are
\[
g_j(F)=\int_p^\infty F(y)[1-G_j(y)]\,dy,\qquad
\gamma(F)=\int_p^\infty F(y)D(y)\,dy.
\]
Subtracting the integrals at two ordered laws gives the stated profit differences. Subtracting the profit-gap
integrals and using the split identity gives the spread difference. Their integrands are nonnegative, so the
monotonicity and exact strictness conditions follow.

Because \(\gamma(r')=\gamma(r)-[\DT(r')-\DT(r)]\), direct subtraction of the threshold ratios gives its displayed
increment. Its numerator is nonnegative when \(\tau(r)\ge0\), and its denominator is positive. Expanding
\(\DT(r')\tau(r')-\DT(r)\tau(r)\) gives the product increment. Both terms are nonnegative, which proves the
product comparison and its exact strictness condition.

On a differentiable path, write \(1-\tau=(g_H-c_H)/\gamma\) and \(\gamma'=-\DT'\). Differentiation gives
\[
J'=\frac{(g_H-c_H)\DT'+\DT g_H'}{\gamma}
       +\frac{\DT(g_H-c_H)\DT'}{\gamma^2}.
\]
Collecting the \(\DT'\) terms and using \(\gamma+\DT=\bar\Gamma\) gives the stated expression.
For the translated uniform laws every incumbent value lies strictly between \(\ell\) and \(h\).
Thus \(t_L=\ell\), \(t_H=\E R=r\), \(g_L=0\) and \(g_H=h-r\), which gives all four formulas and the endpoint
values of \(J\) by substitution. In the uniform-incumbent case, Proposition~\ref{prop:coexistence}(c) gives that
\(\DT(g_H-c_H)\) is nondecreasing. Since \(J=\DT(g_H-c_H)/\gamma\), \(\gamma=\bar\Gamma-\DT\) falls strictly where
\(\DT\) rises, and affordability gives \(g_H>c_H\), \(J\) rises on \([r_0,r_2]\).
\end{proof}

\paragraph{The liquidity share.}\label{par:beta-scope}
The bound \(\beta\le1/3\) is the inequality \(b_0\ge b\). It extends the posterior bounds of the nonzero orders
to the zero order, and through it enters only Corollary~\ref{cor:exclusion} and the zero-order part of the
research bounds. The research parameter region and the slot-feasibility parts of identification range over the
slots with \(\beta\le1/3\). Every other result holds for every \(0<\beta<1/2\); a corollary that fixes the slot
holds at that slot.

\paragraph{A price reached by a nonzero order.}
For every profile and every price function, a price reached by at least one nonzero order has a price posterior
in \([m(\pi),M(\pi)]\), for every \(0<\beta<1/2\), including when the price pools the zero order.

\begin{proof}
Let \(X(\bar P)=\{x:P(x)=\bar P\}\), and suppose it contains a nonzero order. Its liquidity mass
\((1-\lambda)\sum_{x\in X(\bar P)}\eta(x)\) is at least \(b\). Its informed mass in each class lies between zero
and \(\lambda\). Thus the ratio of its class probabilities lies between \((1+u)^{-1}\) and \(1+u\). Bayes' rule
gives the stated interval.
\end{proof}
